\subsection{Joins \& Relational Thinking}

\subsubsection{Why joins exist (relational model intuition)}
\J A \texttt{JOIN} combines rows from two tables based on a matching condition (usually a key 
relationship). Joins let you query related data without storing everything in one table.

\M Conceptually, a join builds a combined row set: it takes rows from table A and table B, 
applies the join condition (for example, \texttt{a.customer\_id = c.id}), and returns matching 
combinations depending on join type.

\Sx Join only what you need (select the required columns). Prefer clear join keys and enforce them 
with constraints when possible. Always validate expected cardinality (one-to-one, one-to-many, 
many-to-many) to prevent row explosions.

\E Performance-wise, the optimizer can choose different join strategies (nested loop, hash join, 
merge join). You don't need internals early, but you must learn to read plans later.

\Pitfall Joining tables without understanding cardinality can duplicate rows. Joining without 
filtering early can create huge intermediate sets.

\subsubsection{INNER JOIN vs LEFT JOIN (and when results change)}
\J \texttt{INNER JOIN} returns only rows that match in both tables. \texttt{LEFT JOIN} returns 
all rows from the left table, plus matching rows from the right; non-matching right columns become 
\texttt{NULL}.

\M \texttt{LEFT JOIN} preserves left-side rows even when no match exists. This is how you query 
''all customers and their orders (if any)''.

\Sx Choose join type based on the business question: ''only those with matches'' $\rightarrow$ inner; 
''include even those without matches'' $\rightarrow$ left. Be careful: placing conditions in 
\texttt{WHERE} can turn a LEFT JOIN into an INNER JOIN unintentionally.

\E In PostgreSQL, the planner can reorder joins, but outer join semantics restrict reordering. 
Learn to spot when the join type forces plan shape.

\Pitfall Using LEFT JOIN but filtering the right table in \texttt{WHERE} drops null-extended rows.

\subsubsection{WHERE vs ON in outer joins (critical senior topic)}
\J \texttt{ON} defines how rows match during the join. \texttt{WHERE} filters rows after the join.

\M With \texttt{LEFT JOIN}, filtering conditions on the right table belong in \texttt{ON} if you 
still want unmatched left rows.

\Sx If you want to preserve left rows, put right-table conditions in \texttt{ON}. If you want 
to drop unmatched rows, put conditions in \texttt{WHERE} (effectively inner join behavior).

\E This is one of the most common causes of ''why did my LEFT JOIN act like INNER JOIN?'' bugs in 
production.

\Pitfall Right-table filters in \texttt{WHERE} eliminate NULL-extended rows and change the result set.

\subsubsection{Join multiplicity (why row counts ''explode'')}
\J If one row in table A matches multiple rows in table B, the result will contain multiple rows 
for that A row.

\M One-to-one: 1 match $\rightarrow$ 1 output row. One-to-many: 1 match $\rightarrow$ many output 
rows. Many-to-many: many matches on both sides $\rightarrow$ potentially huge outputs.

\Sx Always predict expected cardinality before joining. If you need one row per entity, consider 
aggregation (\texttt{GROUP BY}), window functions (for example \texttt{ROW\_NUMBER} to select 
''latest''), and use \texttt{DISTINCT} with caution.

\E Duplicate explosion is often a symptom of poor join conditions or missing constraints 
(for example, missing unique keys).

\Pitfall Using \texttt{DISTINCT} to ''fix'' duplicates without understanding why they happen, hides 
bugs and can be expensive.

\subsubsection{Anti-joins: finding ''missing'' relationships (NOT EXISTS)}
\J An anti-join finds rows in A that have no match in B.

\M The safest pattern is \texttt{NOT EXISTS (subquery)}. This avoids common NULL issues seen with 
\texttt{NOT IN}.

\Sx Prefer \texttt{NOT EXISTS} over \texttt{NOT IN} when the subquery might return \texttt{NULL}. 
Keep anti-join subqueries selective and indexed where possible.

\E PostgreSQL can optimize \texttt{NOT EXISTS} efficiently in many cases, especially with proper 
indexes.

\Pitfall Using \texttt{NOT IN (subquery)} can return empty results if the subquery contains NULL.
